
    %%%%%%%%%%%%%%%%%%%%%%%%%%%%%%%%%%%%%%%%%%%%%%%%%%%%%%%%%%%%%%%%%%%%%%%%%%%%%%%%
    %%% Classe do documento
    \documentclass[12pt, addpoints]{exam}
    \UseRawInputEncoding

    %%%%%%%%%%%%%%%%%%%%%%%%%%%%%%%%%%%%%%%%%%%%%%%%%%%%%%%%%%%%%%%%%%%%%%%%%%%%%%%%
    % preambulo.tex
    %
    % Arquivo de configuração de packages para uso o LaTeX, conforme minhas
    % preferências e modelos pessoais.
    %
    % Para maiores informações, visite:
    %    https://github.com/abrantesasf/latex
    %
    % NÃO ALTERE SE NÃO SOUBER O QUE ESTÁ FAZENDO!
    %%%%%%%%%%%%%%%%%%%%%%%%%%%%%%%%%%%%%%%%%%%%%%%%%%%%%%%%%%%%%%%%%%%%%%%%%%%%%%%%


    %%%%%%%%%%%%%%%%%%%%%%%%%%%%%%%%%%%%%%%%%%%%%%%%%%%%%%%%%%%%%%%%%%%%%%%%%%%%%%%%
    %%% Estruturas de controle:
    \input{static/utils/controles}

    %%%%%%%%%%%%%%%%%%%%%%%%%%%%%%%%%%%%%%%%%%%%%%%%%%%%%%%%%%%%%%%%%%%%%%%%%%%%%%%%
    %%% Compilação condicional em PDF:
    \input{static/utils/pdf}

    %%%%%%%%%%%%%%%%%%%%%%%%%%%%%%%%%%%%%%%%%%%%%%%%%%%%%%%%%%%%%%%%%%%%%%%%%%%%%%%%
    %%% Configurações de layout da página:
    \input{static/utils/layout}

    %%%%%%%%%%%%%%%%%%%%%%%%%%%%%%%%%%%%%%%%%%%%%%%%%%%%%%%%%%%%%%%%%%%%%%%%%%%%%%%%
    %%% Configurações para autores:
    \input{static/utils/autores}

    %%%%%%%%%%%%%%%%%%%%%%%%%%%%%%%%%%%%%%%%%%%%%%%%%%%%%%%%%%%%%%%%%%%%%%%%%%%%%%%%
    %%% Cabeçalho e rodapé:
    %%% Atenção! É MUITO DIFÍCIL ajustar apropriadamente os running headers
    %%% e running footers da primeira vez. Você pode confiar no LaTeX e deixar que
    %%% ele faça o ajuste automaticamente ou pode quebrar a cabeça e
    %%% tentar melhorar algo que o LaTeX já faz muito bem, mesmo que não fique
    %%% exatamente do jeito que você quer. O que você prefere? Se quiser ajustar
    %%% manualmente, descomente a linha abaixo e faça as configurações no arquivo
    %%% "utils/headings.tex".
    %\input{static/utils/headings}

    %%%%%%%%%%%%%%%%%%%%%%%%%%%%%%%%%%%%%%%%%%%%%%%%%%%%%%%%%%%%%%%%%%%%%%%%%%%%%%%%
    %%% Configuração para as fontes do documento:
    %%% Se você quiser usar fontes próprias ou do sistema, precisará ajustar as
    %%% configurações no arquivo "utils/fontes.tex".
    %%% ATENÇÃO: por padrão eu utilizo XeLaTeX com fontes proprietárias projetadas
    %%% por Matthew Butterick, adquiridas em https://mbtype.com, que não podem ser
    %%% distribuídas devido à licença de utilização. Se você usar XeLaTeX, você
    %%% DEVERÁ OBRIGATORIAMENTE ajustar as fontes no arquivo "utils/fontes.tex".
    \input{static/utils/fontes}

    %%%%%%%%%%%%%%%%%%%%%%%%%%%%%%%%%%%%%%%%%%%%%%%%%%%%%%%%%%%%%%%%%%%%%%%%%%%%%%%%
    %%% Sumário:
    \input{static/utils/toc}

    %%%%%%%%%%%%%%%%%%%%%%%%%%%%%%%%%%%%%%%%%%%%%%%%%%%%%%%%%%%%%%%%%%%%%%%%%%%%%%%%
    %%% Matemática:
    \input{static/utils/matematica}

    %%%%%%%%%%%%%%%%%%%%%%%%%%%%%%%%%%%%%%%%%%%%%%%%%%%%%%%%%%%%%%%%%%%%%%%%%%%%%%%%
    %%% Notas de rodapé:
    \input{static/utils/footnotes}

    %%%%%%%%%%%%%%%%%%%%%%%%%%%%%%%%%%%%%%%%%%%%%%%%%%%%%%%%%%%%%%%%%%%%%%%%%%%%%%%%
    %%% Referências cruzadas, links, citações:
    \input{static/utils/crossrefs}

    %%%%%%%%%%%%%%%%%%%%%%%%%%%%%%%%%%%%%%%%%%%%%%%%%%%%%%%%%%%%%%%%%%%%%%%%%%%%%%%%
    %%% Configurações para os hiperlinks em PDF:
    \input{static/utils/pdfconfig}

    %%%%%%%%%%%%%%%%%%%%%%%%%%%%%%%%%%%%%%%%%%%%%%%%%%%%%%%%%%%%%%%%%%%%%%%%%%%%%%%%
    %%% Glossário e índice remissivo:
    \input{static/utils/glossindex}

    %%%%%%%%%%%%%%%%%%%%%%%%%%%%%%%%%%%%%%%%%%%%%%%%%%%%%%%%%%%%%%%%%%%%%%%%%%%%%%%%
    %%% Referências bibliográficas:
    \input{static/utils/bibliografia}

    %%%%%%%%%%%%%%%%%%%%%%%%%%%%%%%%%%%%%%%%%%%%%%%%%%%%%%%%%%%%%%%%%%%%%%%%%%%%%%%%
    %%% Cores:
    \input{static/utils/cores}

    %%%%%%%%%%%%%%%%%%%%%%%%%%%%%%%%%%%%%%%%%%%%%%%%%%%%%%%%%%%%%%%%%%%%%%%%%%%%%%%%
    %%% Computação:
    \input{static/utils/computacao}

    %%%%%%%%%%%%%%%%%%%%%%%%%%%%%%%%%%%%%%%%%%%%%%%%%%%%%%%%%%%%%%%%%%%%%%%%%%%%%%%%
    %%% Gráficos:
    \input{static/utils/graficos}

    %%%%%%%%%%%%%%%%%%%%%%%%%%%%%%%%%%%%%%%%%%%%%%%%%%%%%%%%%%%%%%%%%%%%%%%%%%%%%%%%
    %%% Ícones:
    \input{static/utils/icones}

    %%%%%%%%%%%%%%%%%%%%%%%%%%%%%%%%%%%%%%%%%%%%%%%%%%%%%%%%%%%%%%%%%%%%%%%%%%%%%%%%
    %%% Blocos:
    \input{static/utils/blocos}

    %%%%%%%%%%%%%%%%%%%%%%%%%%%%%%%%%%%%%%%%%%%%%%%%%%%%%%%%%%%%%%%%%%%%%%%%%%%%%%%%
    %%% Boxes coloridos:
    \input{static/utils/colorbox}

    %%%%%%%%%%%%%%%%%%%%%%%%%%%%%%%%%%%%%%%%%%%%%%%%%%%%%%%%%%%%%%%%%%%%%%%%%%%%%%%%
    %%% Tabelas:
    \input{static/utils/tabelas}

    %%%%%%%%%%%%%%%%%%%%%%%%%%%%%%%%%%%%%%%%%%%%%%%%%%%%%%%%%%%%%%%%%%%%%%%%%%%%%%%%
    %%% Ambientes de listas:
    \input{static/utils/listas}

    %%%%%%%%%%%%%%%%%%%%%%%%%%%%%%%%%%%%%%%%%%%%%%%%%%%%%%%%%%%%%%%%%%%%%%%%%%%%%%%%
    %%% Ambientes floats e similares:
    \input{static/utils/floats}

    %%%%%%%%%%%%%%%%%%%%%%%%%%%%%%%%%%%%%%%%%%%%%%%%%%%%%%%%%%%%%%%%%%%%%%%%%%%%%%%%
    %%% Ajustes para a classe EXAM:
    \input{static/utils/exam}

    %%%%%%%%%%%%%%%%%%%%%%%%%%%%%%%%%%%%%%%%%%%%%%%%%%%%%%%%%%%%%%%%%%%%%%%%%%%%%%%%
    %%% Ajustes para textos:
    \input{static/utils/texto}

    %%%%%%%%%%%%%%%%%%%%%%%%%%%%%%%%%%%%%%%%%%%%%%%%%%%%%%%%%%%%%%%%%%%%%%%%%%%%%%%%
    %%% Ambientes, aliases, comandos e customizações em geral para serem
    %%% utilizadas nos documentos. Defina aqui qualquer customização específica
    %%% para seus documentos!
    \input{static/utils/customizacoes}

    %%%%%%%%%%%%%%%%%%%%%%%%%%%%%%%%%%%%%%%%%%%%%%%%%%%%%%%%%%%%%%%%%%%%%%%%%%%%%%%%
    %%% Ajuste do layout, espaçamento de linhas e etc.:
    \geometry{a4paper, portrait, left=2.0cm, right=2.0cm, top=2.0cm, bottom=2.0cm}
    \singlespacing

    %%%%%%%%%%%%%%%%%%%%%%%%%%%%%%%%%%%%%%%%%%%%%%%%%%%%%%%%%%%%%%%%%%%%%%%%%%%%%%%%
    %%% Configurações de cabeçalho e rodapé (não use fancyhdr pois ocorre conflito):
    \pagestyle{headandfoot}
    \runningheadrule
    \firstpageheader{}{}{}
    \runningheader{Sem disciplina}{}{Avaliação}
    \firstpagefooter{}{}{Página \thepage\ de \numpages}
    \runningfooter{}{}{Página \thepage\ de \numpages}
    \runningfootrule

    %%%%%%%%%%%%%%%%%%%%%%%%%%%%%%%%%%%%%%%%%%%%%%%%%%%%%%%%%%%%%%%%%%%%%%%%%%%%%%%%
    %%% Configurações para as propriedades do PDF:
    \hypersetup{
    hidelinks,           % Comente para web, descomente para impressão
    colorlinks=true,      % True para web, False para impressão
    pdftitle=tesste: Avaliação,
    pdfauthor=b,
    pdfsubject=Sem disciplina,
    pdfkeywords=['Sem tag'],
    pdfinfo={
        CreationDate={}, % Ex.: D:AAAAMMDDHH24MISS
        ModDate={}       % Ex.: D:AAAAMMDDHH24MISS
    }
    }
    \input{static/utils/pdfconfig}

    %%%%%%%%%%%%%%%%%%%%%%%%%%%%%%%%%%%%%%%%%%%%%%%%%%%%%%%%%%%%%%%%%%%%%%%%%%%%%%%%
    %%% Começa o documento
    \begin{document}

    %%%%%%%%%%%%%%%%%%%%%%%%%%%%%%%%%%%%%%%%%%%%%%%%%%%%%%%%%%%%%%%%%%%%%%%%%%%%%%%%
    %%% Folha de instruções padronizadas da UVV para a prova
    \pdfbookmark[1]{Instruções}{instrucoes}
    \begin{figure}[H]
        \begin{center}
            \includegraphics[scale=0.3]{{static/imgs/logo-uvv.png}}
        \end{center}	
    \end{figure}

    \vspace{-1cm}

    \begin{table}[H]
        \centering
        \begin{tabular}{|llll|}
            \hline
            \multicolumn{2}{|l|}{\textbf{Disciplina:} Sem disciplina} & 
            \multicolumn{1}{l|}{\multirow{3}{*}{\textbf{\makecell{Nota:\\ \,\\ \,}}}} & 
            \multirow{3}{*}{\textbf{\makecell{Coordenador:\\ \,\\ \,}}} \\ 
            \cline{1-2}
            \multicolumn{2}{|l|}{\textbf{Professor:} b \hspace{4.72cm} } & 
            \multicolumn{1}{l|}{} & \\ 
            \cline{1-2}
            \multicolumn{2}{|l|}{\textbf{Aluno:}} & 
            \multicolumn{1}{l|}{} & \\ 
            \hline
            \multicolumn{1}{|l|}{\textbf{Turma:} \hspace{6.3cm} } & 
            \multicolumn{1}{l|}{\textbf{Semestre:}} & 
            \multicolumn{2}{l|}{\textbf{Valor:} 10 pontos} \\ 
            \hline
            \multicolumn{1}{|l|}{\textbf{Data:}} & 
            \multicolumn{3}{l|}{\textbf{Avaliação:}} \\ 
            \hline
        \end{tabular}
    \end{table}

    \begin{center}
        \fbox{\setlength\fboxsep{0.3cm} \fbox{\parbox{15.4cm}{
            \begin{center}
                \textbf{INSTRUÇÕES PARA A REALIZAÇÃO DESTA AVALIAÇÃO:}
            \end{center}
            \begin{itemize}[leftmargin=*]
                \item Esta é a primeira avaliação da disciplina Sem disciplina, 
                    e engloba o conteúdo referente Sem tag.
                \item Esta avaliação contém 10 questões objetivas, \textbf{todas obrigatórias}.
                \item \textbf{Leia atentamente cada questão!} As questões objetivas terão uma,
                    e apenas uma única, resposta a ser marcada.
                \item \textbf{Desligue o celular} antes de começar. A avaliação é
                    \textbf{individual} e \textbf{sem consulta}.
                \item As questões podem ser respondidas, na prova, com lápis ou caneta. Ao final
                    da prova \textbf{{VOCÊ DEVERÁ PREENCHER O GABARITO COM CANETA
                    DE COR PRETA}} \textbf{\underline{OU AZUL}} para que seja feita a correção
                    automatizada através da leitura do padrão de respostas marcado no
                    gabarito.
                \item \textbf{CUIDADO ao preencher o gabarito} pois eles são \textbf{INDIVIDUAIS
                    e NOMEADOS} (cada estudante tem um gabarito próprio específico, com um
                    código de identificação único). \textbf{Questões rasuradas são anuladas}
                    pelo software de leitura do gabarito.
                \item Ao final da prova, \textbf{DEVOLVA A PROVA E O GABARITO} para o professor.
                \item Siga todas as normas de \textbf{Integridade Acadêmica} da disciplina.
                    Alunos que forem flagrados com qualquer espécie de ``cola'' ou trocando
                    informações com outros alunos terão suas avaliações recolhidas, as notas
                    zeradas, e a situação será encaminhada para a coordenação para a aplicação
                    das penalidades previstas pela universidade.
                \item Com 90 minutos de avaliação você terá tempo de até 9,min por questão,
                    em média.
                \item Boa avaliação!
            \end{itemize}
            \vspace{2cm}
        }}}
    \end{center}
    \newpage

    %%%%%%%%%%%%%%%%%%%%%%%%%%%%%%%%%%%%%%%%%%%%%%%%%%%%%%%%%%%%%%%%%%%%%%%%%%%%%%%%
    %%% Inicia o bloco de questões:
    \begin{questions}

    \newpage
    \question
No que diz respeito as vantagens da arquitetura de micro-núcleo para sistemas operacionais em relação a arquiteturas de núcleo monolítico, quais das seguintes afirmações são verdadeiras?

\begin{itemize}
    \item[I] A arquitetura de micro-núcleo facilita a depuração do SO.
    \item[II] A arquitetura de micro-núcleo permite um número menor de mudanças de contexto.
    \item[III] A arquitetura de micro-núcleo facilita a reconfiguração de serviços do SO pois a maioria deles reside em espaço de usuário.
\end{itemize}
\begin{choices}
\choice I e III
\choice II e III
\choice Todas são verdadeiras
\choice I e II
\choice Apenas I
\end{choices}

\question
Pode-se afirmar que o gráfico da função \(y = 2 + \frac{1}{x - 1}\) é o gráfico da função \(y = \frac{1}{x}\):
\begin{choices}
\choice transladado uma unidade para a direita e duas unidades para baixo
\choice transladado uma unidade para a esquerda e duas unidades para cima
\choice transladado uma unidade para a direita e duas unidades para cima
\choice nenhuma das anteriores.
\choice transladado uma unidade para a esquerda e duas unidades para baixo
\end{choices}

\question
Supondo a Relação $\texttt{PROJ (PNO, Orçam)}$, com chave primária $\texttt{PNO}$, a Relação $\texttt{EMP (ENO, ENome, Cargo)}$ com chave primária $\texttt{ENO}$, e a Relação $\texttt{DSG (ENO, PNO, Dur, Resp)}$, com chave primária \{\texttt{ENO, PNO}\}, chave estrangeira $\texttt{PNO}$ em relação a $\texttt{PROJ}$ e chave estrangeira $\texttt{ENO}$ em relação a $\texttt{EMP}$. Qual das expressões da álgebra relacional abaixo $\textbf{NÃO}$ corresponde à seguinte consulta SQL:

\begin{verbatim}
SELECT ENome
FROM EMP, PROJ, DSG
WHERE EMP.ENO = DSG.ENO
      AND PROJ.PNO = DSG.PNO
      AND Dur > 36
\end{verbatim}
\begin{choices}
\choice \( \pi_{ENome} (\text{PROJ} \Join_{PNO} (\text{EMP} \Join_{ENO} \sigma_{Dur > 36} (\pi_{Dur} (\text{DSG})))) \)
\choice \( \pi_{ENome} (\sigma_{Dur > 36} ((\pi_{ENome, ENO} (\text{EMP})) \Join_{ENO} (\text{PROJ} \Join_{PNO} (\text{DSG})))) \)
\choice \( \pi_{ENome} (\text{PROJ} \Join_{PNO} (\text{EMP} \Join_{ENO} \sigma_{Dur > 36} (\text{DSG}))) \)
\choice \( \pi_{ENome} (\text{PROJ} \Join_{PNO} (\sigma_{Dur > 36} (\text{EMP} \Join_{ENO} (\text{DSG})))) \)
\choice \( \pi_{ENome} (\sigma_{Dur > 36} ((\pi_{PNO} (\text{PROJ})) \Join_{PNO} (\text{EMP} \Join_{ENO} (\pi_{Dur} (\text{DSG})))) \)
\end{choices}

\question
Em relação ao paradigma de programação cliente-servidor. Qual das afirmativas abaixo é FALSA?
\begin{choices}
\choice Um aplicativo cliente pode acessar múltiplos serviços quando necessário.
\choice Um aplicativo servidor inicia ativamente o contato com clientes 
\choice Um aplicativo servidor aceita contato de clientes arbitrários, mas oferece um único serviço. 
\choice Um aplicativo cliente é um programa arbitrário que se torna temporariamente um cliente quando for necessário o acesso remoto a um serviço, mas também executa processamento local.
\choice Um aplicativo servidor é um programa de propósito especial dedicado a fornecer um serviço, mas pode tratar de múltiplos clientes remotos ao mesmo tempo. 
\end{choices}

\question
Considere a seguinte tabela em uma base de dados relacional (chave primária sublinhada):\\
Tabela1 (CodAluno, CodDisciplina, AnoSemestre, NomeAluno, NomeDisciplina, CodNota, DescricaoNota)\\
Considere as seguintes dependências funcionais:
\begin{itemize}
    \item CodAluno \(\rightarrow\) NomeAluno
    \item CodDisciplina \(\rightarrow\) NomeDisciplina
    \item (CodAluno, CodDisciplina, AnoSemestre) \(\rightarrow\) CodNota
    \item (CodAluno, CodDisciplina, AnoSemestre) \(\rightarrow\) DescricaoNota
    \item CodNota \(\rightarrow\) DescricaoNota
\end{itemize}

Considerando as formas normais, qual das afirmativas abaixo se aplica:
\begin{choices}
\choice A tabela encontra-se na terceira forma normal, mas não na quarta forma normal.
\choice A tabela encontra-se na primeira forma normal, mas não na segunda forma normal.
\choice A tabela encontra-se na segunda forma normal, mas não na terceira forma normal.
\choice A tabela não está na primeira forma normal.
\choice A tabela está na quarta forma normal.
\end{choices}

\question
Considere as seguintes tabelas em uma base de dados relacional:\\
Departamento (\underline{CodDepto}, NomeDepto)\\
Empregado (\underline{CodEmp}, NomeEmp, CodDepto, Salario)\\
Considere a seguinte consulta SQL:
\begin{verbatim}
SELECT d.CodDepto, d.NomeDepto, SUM(e.Salario)
FROM Departamento d
JOIN Empregado e ON d.CodDepto = e.CodDepto
GROUP BY d.CodDepto, d.NomeDepto
HAVING COUNT(e.CodEmp) > 2 AND AVG(e.Salario) > 40;
\end{verbatim}

Qual o resultado da consulta?
\begin{choices}
\choice Para cada departamento que tem mais que dois empregados e cuja média salarial é maior que 40, obter o código de departamento, seguido do nome do departamento, seguido da soma dos salários dos empregados do departamento.
\choice Para cada departamento que tem mais que dois empregados e cuja média salarial é maior que 40, obter um grupo de linhas que contém, para cada empregado do departamento, o código de seu departamento, seguido do nome de seu departamento, seguido da soma dos salários dos empregados do departamento.
\choice A consulta não retorna nada pois está incorreta.
\choice Para cada empregado que tem mais que dois departamentos, ambos com média salarial maior que 40, obter o código de departamento, seguido do nome do departamento, seguido da soma dos salários dos empregados do departamento.
\choice Para cada departamento que tem mais que dois empregados e cuja média salarial, considerando todos empregados do departamento, exceto os dois primeiros, é maior que 40, obter o código de departamento, seguido do nome do departamento, seguido da soma dos salários dos empregados do departamento.
\end{choices}

\question
(ENADE 2017) No modelo relacional de dados, uma junção entre tabelas cria uma
outra tabela derivada de duas ou mais tabelas de acordo com os critérios
especificados para a junção, de modo similar às regras da teoria dos
conjuntos. Nos conjuntos a seguir, as áreas hachuradas representam os resultados
das consultas SQL em que ``id' é uma chave primária e ``fk' é uma chave
estrangeira:
\begin{figure}[H]
  \centering
  \includegraphics[width=0.5\linewidth]{static/imgs/defaul.png}
\end{figure}
Assinale a opção em que a declaração SQL é equivalente ao conjunto apresentado a
seguir:
\begin{figure}[H]
  \centering
  \includegraphics[width=0.5\linewidth]{static/imgs/defaul.png}
\end{figure}

\begin{choices}
\choice \begin{verbatim}SELECT *
FROM tabela1 t1
RIGHT JOIN tabela2 t2 ON (t1.id = t2.fk)
WHERE t2.fk <> NULL;
\end{verbatim}

\choice \begin{verbatim}SELECT *
FROM tabela1 t1 WHERE EXISTS (SELECT 1
                              FROM tabela2 t2
                              WHERE t2.id = t1.fk);
\end{verbatim}

\choice \begin{verbatim}SELECT *
FROM tabela1 t1
INNER JOIN tabela2 t2 ON (t1.id = t2.fk);
\end{verbatim}

\choice \begin{verbatim}SELECT *
FROM tabela1 t1
WHERE NOT EXISTS (SELECT 1
                  FROM tabela1 t2
                  WHERE t2.id = tk.fk);
\end{verbatim}

\choice \begin{verbatim}SELECT *
FROM tabela1 t1
RIGHT JOIN tabela2 t2 ON (t1.id = t2.fk)
WHERE t1.id IS NULL;
\end{verbatim}

\end{choices}

\question
Uma característica de uma arquitetura RISC é:
\begin{choices}
\choice Possui um pequeno conjunto de instruções simples
\choice A arquitetura é constituída de milhares de processadores
\choice Uma arquitetura de alto risco pois o mercado de hardware evolui muito rapidamente
\choice O processador é formado por válvulas e transistores
\choice Possui um grande conjunto de instruções complexas
\end{choices}

\question
Quando trabalhando com sistemas baseados em trocas de mensagens, temporizações (\textit{time-outs}) são utilizadas para:

\begin{choices}
\choice Limitar o tamanho de uma mensagem transmitida.
\choice Temporariamente suspender a transmissão de mensagens.
\choice Limitar o tempo para obter um recurso.
\choice Arbitrar que uma mensagem transmitida foi perdida.
\choice Limitar o número de retransmissões de uma mensagem.
\end{choices}

\question
Um sistema centralizado é um concentrador de recursos; um sistema distribuído apresenta seus recursos dispersos. Entretanto, nem todo conjunto de recursos computacionais dispersos pode ser considerado um sistema distribuído. Considerando um conjunto de computadores, assinale a alternativa que melhor corresponde às características necessárias para considerá-lo um sistema distribuído:
\begin{choices}
\choice Existência de sistema operacional idêntico e hardware padronizado em todos os computadores
\choice Suporte de rede e funções primitivas de comunicação 
\choice Existência de memória compartilhada e relógios locais sincronizados 
\choice Suporte de rede e um relógio global
\choice Existência de memória secundária compartilhada e protocolos de sincronização de estado
\end{choices}



    %%%%%%%%%%%%%%%%%%%%%%%%%%%%%%%%%%%%%%%%%%%%%%%%%%%%%%%%%%%%%%%%%%%%%%%%%%%%%%%%
    %%% Finaliza o bloco de questões:
    \end{questions}
    \end{document}
    